\subsection{引言}\label{sec:12.1}
\autoref{sec:11}讨论了可称为“经典线性模型”的内容. 方差分析和简单线性回归都基于具有正态误差的线性模型. 本章考察这一模型的一些推广，它们在实际问题中已证明很有用.

在\autoref{sec:12.2}中，将线性模型推广到预测变量含误差的模型，称为变量含误差的回归（regression with errors in variables，EIV）. 此时预测变量$X$像响应变量$Y$一样成为随机变量. 这一模型中的估计会遇到许多意想不到的困难，与简单线性回归模型可能很不相同.

\autoref{sec:12.3}进一步推广线性模型，考察逻辑回归（logistic regression）. 此时响应变量是离散的Bernoulli变量. Bernoulli均值是有界函数，在有界函数上建立线性模型可能遇到问题，尤其在边界附近. 因此，用logit变换将均值变为无界参数，再把变换后的参数建模为预测变量的线性函数. 当对响应均值的某个函数建立线性模型时，就得到广义线性模型（generalized linear model）.

最后，在\hyperref[sec:12.4]{节11.3}中考察线性回归中的稳健性.\footnote{译注：原文这里写“Section 11.3”，按本章内容应为\autoref{sec:12.4}. 原编号保留.} 本章其他节改变模型，这里则改变拟合准则. 其发展与\autoref{sec:10.2.2}关于稳健点估计的讨论平行. 即以基于$\rho$函数的准则代替最小二乘准则，使估计对观测值较不敏感，同时保留一定效率.
